\documentclass{ip-quiz}

\quizlevel{n2}
\quiznum{2}
\quiztitle{Strings y condicionales}

\begin{document}

\makeheader
\maketitle

\begin{sectionbox}{Enunciado}
Un estudiante busca en su computador el ensayo final que debe enviar en contados minutos, entre los múltiples archivos de borrador que generó en su afán.

Escriba una función que recibe cuatro nombres de archivo y determine cuál es el ensayo final. Un archivo es el ensayo final si: su nombre contiene la palabra \texttt{ensayo}, está marcado como la tercera versión con el indicador \texttt{v3}, es un archivo PDF y la palabra \texttt{final} aparece tres veces.

Si exactamente uno cumple todas las condiciones, retórnelo. De lo contrario, priorice el primero en orden alfabético. Si persiste el empate, retorne el primero en el orden de los parámetros.

Los nombres de los archivos no tienen tildes.
\end{sectionbox}

\pagebreak

\begin{sectionbox}{Solución}
\begin{minted}[escapeinside=||]{python}
def cumple_condiciones(nombre_ensayo: str) -> bool:
    |\tikzmark{cctop}|return ("ensayo" in nombre_ensayo and|\tikzmark{ccright}|
            "v3" in nombre_ensayo and
            nombre_ensayo[-4:] == ".pdf" and
            nombre_ensayo.count("final") == 3)|\tikzmark{ccbot}|


def es_mejor_candidato(
    candidato: str,
    mejor: str,
    ningun_ensayo_cumple: bool
) -> bool:
    |\tikzmark{nincmark}||\tms{algA}|if ningun_ensayo_cumple:|\tme{algA}|
        return candidato.lower() < mejor.lower()
    return (cumple_condiciones(candidato) and
            (not cumple_condiciones(mejor) or |\tms{lwrA}|candidato.lower() < mejor.lower()|\tme{lwrA}|))


def encontrar_ensayo_final(
    ensayo_1: str, 
    ensayo_2: str, 
    ensayo_3: str, 
    ensayo_4: str
) -> str:
    cumple_1 = cumple_condiciones(ensayo_1)
    cumple_2 = cumple_condiciones(ensayo_2)
    cumple_3 = cumple_condiciones(ensayo_3)
    cumple_4 = cumple_condiciones(ensayo_4)

    alguno_cumple = cumple_1 or cumple_2 or cumple_3 or cumple_4

    mejor = ensayo_1
    if |\tms{ecA}|es_mejor_candidato|\tme{ecA}|(ensayo_2, mejor, not alguno_cumple):
        mejor = ensayo_2
    if es_mejor_candidato(ensayo_3, mejor, not alguno_cumple):
        mejor = ensayo_3
    if es_mejor_candidato(ensayo_4, mejor, not alguno_cumple):
        mejor = ensayo_4

    return mejor
\end{minted}

\begin{tikzpicture}[remember picture, overlay]

  \codebrace{ccright}{1.4cm}{cctop}{ccbot}{3pt}{5.5cm}
    {Podrían haber sido varios \texttt{if}s, pero como todas las condiciones deben cumplirse, basta con usar \texttt{and}.}

  \node[box, fit=(algAs)(algAe)] (calgA) {};
  \node[note, anchor=north west, text width=5.5cm]
    (noteAlg) at ([xshift=-8.9cm, yshift=0.9cm] current page.east |- nincmark)
    {Contempla el caso en que ningún ensayo cumple con las condiciones: en ese caso la comparación es puramente alfabética, no por condiciones.};
  \draw[arr, dotted] (calgA.east) to[bend right=5] (noteAlg.west);

  \node[box, fit=(ecAs)(ecAe)] (cecA) {};
  \node[note, anchor=north west, text width=4.5cm]
    (noteEc) at ([xshift=-7.5cm, yshift=0.3cm] current page.east |- ecAs)
    {Comparación sucesiva con el mejor candidato. Para no repetir la misma lógica tres veces, la encapsulo en una función.};
  \draw[arr, dotted] (cecA.east) to[bend right=5] (noteEc.west);

  \node[box, fit=(lwrAs)(lwrAe)] (clwrA) {};
  \node[note, anchor=north, text width=4.5cm]
    (noteLwr) at ([xshift=-8.5cm, yshift=-0.5cm] current page.east |- lwrAs)
    {\texttt{.lower()} para que la comparación alfabética sea correcta, porque en Python, mayúsculas \textless{} minúsculas en ASCII.};
  \draw[arr, dotted] (clwrA.south) to[bend right=5] (noteLwr.north);

\end{tikzpicture}
\end{sectionbox}

\end{document}
